% !TEX root = ../../thesis.tex

\section{Regularity}\label{sec:regularity}

By Proposition~\ref{thm:CLS-stationary}, the limit varifold $V$ is stationary; by Theorem~\ref{thm:integrality-mult-one} (no mass cancellation) or Proposition~\ref{prop:integrality-nested} (mass cancellation), $V$ is integral. To promote $V$ to a smooth embedded minimal hypersurface, we apply Wickramasekera's regularity theorem, which requires two further inputs: \emph{stability} and the \emph{$\alpha$-structural hypothesis}. This section verifies both, in both the no-mass-cancellation and mass-cancellation cases.

\begin{definition}[$\alpha$-structural hypothesis]\label{defn:alpha-structural}
For $\alpha \in (0, 1/2)$, a stationary integral varifold $V$ satisfies the \emph{$\alpha$-structural hypothesis} if no singular point
of $V$ has a neighborhood in which $V$ corresponds to a union of embedded $C^{1,\alpha}$ hypersurfaces-with-boundary
meeting only along their common $C^{1,\alpha}$ boundary, with multiplicity a constant positive integer
on each of the constituent hypersurfaces-with-boundary.
\end{definition}

\begin{theorem}[Wickramasekera~{\cite{Wic14}}]\label{thm:wickramasekera}
    Let $V$ be a stable integral $n$-varifold in a closed $(n+1)$-dimensional Riemannian manifold
    with $2\leq n\leq 6$. If $V$ satisfies the $\alpha$-structural hypothesis for some $\alpha \in (0, 1/2)$,
    then $\sing V = \emptyset$, i.e., $\spt\|V\|$ is a smooth embedded minimal hypersurface.
\end{theorem}

%%%%%%%%%%%%%%%%%%%%%%%%%%%%%%%%%%%%%%%%%%%%%%%%%%%%%%%%%%%%%%%%%%%%%%%%%%%%%%%
\subsection{Stability}\label{subsec:stability}
%%%%%%%%%%%%%%%%%%%%%%%%%%%%%%%%%%%%%%%%%%%%%%%%%%%%%%%%%%%%%%%%%%%%%%%%%%%%%%%

The non-excessive condition provides stability through the following chain of implications: non-excessiveness ensures that the limit varifold $V$ is one-sided homotopic minimizing at all but finitely many points, and one-sided minimization implies stability.

\begin{proposition}[Stability of the limit varifold]\label{prop:stability}
Let $\Phi$ be a non-excessive ONVP sweepout, $x_i \to x_0 \in \mathfrak{m}(\Phi)$ a min-max sequence, and $V$ the varifold limit of $|\partial^*\Omega(x_i)|$. Then $V$ is stable.
\end{proposition}

\begin{proof}
By~\cite[Proposition~3.1]{CLS22}, the set $\mathfrak{h}_{\mathrm{nm}}(V)$ of non-homotopic-minimizing points (Definition~\ref{def:hnm}) is finite. Therefore $V$ is one-sided homotopic minimizing in $B_r(P)$ for every $P \in \spt\|V\| \setminus \mathfrak{h}_{\mathrm{nm}}(V)$ and some $r = r(P) > 0$. One-sided homotopic minimization implies that the second variation $\delta^2 V(\varphi, \varphi) \geq 0$ for all normal deformations $\varphi$ supported in $B_r(P)$. Since the finitely many points of $\mathfrak{h}_{\mathrm{nm}}(V)$ form a set of $\mathcal{H}^n$-measure zero, a partition-of-unity argument extends stability to all compactly supported normal deformations, giving global stability. Combined with the index bound $\mathrm{Index}(V) \leq 1$ from the min-max construction, $V$ is stable.

We remark that the finiteness of $\mathfrak{h}_{\mathrm{nm}}(V)$ in~\cite[Proposition~3.1]{CLS22} relies on the non-excessive property of the sweepout slices $\partial^*\Omega(x_i)$, not on the identification $V = |\partial^*\Omega(x_0)|$. In particular, the argument applies equally in the mass-cancellation case, where $V \neq |\partial^*\Omega(x_0)|$: the non-excessive condition controls the almost-minimizing behavior of the approximating slices, and the one-sided minimizing property transfers to $V$ through varifold convergence.
\end{proof}

%%%%%%%%%%%%%%%%%%%%%%%%%%%%%%%%%%%%%%%%%%%%%%%%%%%%%%%%%%%%%%%%%%%%%%%%%%%%%%%
\subsection{Tangent cone classification}\label{subsec:tangent-cones}
%%%%%%%%%%%%%%%%%%%%%%%%%%%%%%%%%%%%%%%%%%%%%%%%%%%%%%%%%%%%%%%%%%%%%%%%%%%%%%%

The structure of tangent cones at singular points of $V$ depends on whether mass cancellation occurs.

\medskip\noindent\textbf{No mass cancellation.}
When no mass cancellation occurs, $V = |\partial^*\Omega(x_0)|$ for a Caccioppoli set $\Omega(x_0)$ (Theorem~\ref{thm:integrality-mult-one}). Tangent cones at singular points are constrained by the Caccioppoli set structure. For stationary boundaries of Caccioppoli sets in $\mathbb{R}^{n+1}$ ($n \leq 6$), the only possible singular cones are:
\begin{itemize}
    \item $C = |H_1| + |H_2|$ (two half-hyperplanes meeting at angle $< \pi$) --- a ``wedge'';
    \item $C = 2|H|$ (single half-hyperplane with multiplicity 2) --- a ``fold''.
\end{itemize}
Higher multiplicities or configurations where three or more sheets meet cannot arise, since the boundary of a Caccioppoli set separates $\mathbb{R}^{n+1}$ into exactly two regions ($\Omega$ and $\Omega^c$).

\medskip\noindent\textbf{Mass cancellation.}
When mass cancellation occurs, $V$ is integral (Proposition~\ref{prop:integrality-nested}) but is not \emph{a priori} the boundary of a single Caccioppoli set. Tangent cones at singular points are still stationary integral cones in $\mathbb{R}^{n+1}$, but the two-region constraint no longer applies. In addition to wedges and folds, junctions of $N \geq 3$ half-hyperplanes meeting along a common $(n-1)$-dimensional edge can occur. The $\alpha$-structural hypothesis, if it fails, produces exactly such a junction structure at some $Z \in \sing V$: finitely many $C^{1,\alpha}$ hypersurfaces-with-boundary meeting along a common $C^{1,\alpha}$ boundary $\Gamma$ of dimension $n - 1$.

In both cases, the $\alpha$-structural hypothesis is verified by showing that each possible junction structure admits a local area improvement, leading to a contradiction. The details are carried out in the next subsection.

%%%%%%%%%%%%%%%%%%%%%%%%%%%%%%%%%%%%%%%%%%%%%%%%%%%%%%%%%%%%%%%%%%%%%%%%%%%%%%%
\subsection{Junction resolution and the $\alpha$-structural hypothesis}\label{subsec:junction-resolution}
%%%%%%%%%%%%%%%%%%%%%%%%%%%%%%%%%%%%%%%%%%%%%%%%%%%%%%%%%%%%%%%%%%%%%%%%%%%%%%%

\begin{theorem}[Verification of $\alpha$-structural hypothesis]\label{thm:alpha-hypo}
    Let $V$ be the stationary integral varifold arising from a non-excessive ONVP sweepout $\Phi$
    with width $W$, via the min-max sequence $x_i \to x_0 \in \mathfrak{m}(\Phi)$. Then $V$ satisfies the $\alpha$-structural hypothesis.
\end{theorem}

We treat the no-mass-cancellation and mass-cancellation cases in a unified proof, noting where the arguments diverge.

\begin{proof}
Suppose the $\alpha$-structural hypothesis (Definition~\ref{defn:alpha-structural}) fails at some $Z \in \sing V$. Then there exists $\rho > 0$ such that $\spt\|V\| \cap B_\rho(Z)$ consists of finitely many $C^{1,\alpha}$ embedded hypersurfaces-with-boundary meeting along a common $C^{1,\alpha}$ boundary $\Gamma$ of dimension $n - 1$.

\medskip\noindent\textbf{Step~1: Junction resolution at the tangent cone level.}
Blow up $V$ at $Z$: any tangent cone $C$ is a stationary integral cone in $\mathbb{R}^{n+1}$. The $\alpha$-structural failure gives $C$ a junction structure. By the tangent cone classification (Section~\ref{subsec:tangent-cones}), three cases arise:

\begin{itemize}
    \item \textbf{Case 1a (wedge):} The tangent cone is $C = |H_1| + |H_2|$, where $H_1, H_2$ are half-hyperplanes meeting along their common boundary at angle $\theta < \pi$.

    In the no-mass-cancellation case, write $C = |\bdy\Omega_C|$ for a Caccioppoli set $\Omega_C \subset \mathbb{R}^{n+1}$. Since the angle between $H_1$ and $H_2$ is less than $\pi$, the wedge $\partial\Omega_C$ is not perimeter-minimizing: in each two-dimensional cross-section normal to the edge, the two half-lines of total length $2r$ can be replaced by the chord of length $2r\sin(\theta/2) < 2r$, yielding a Caccioppoli set $\Omega_C' \subset \Omega_C$ with $\Omega_C \triangle \Omega_C' \subset B_1$ and
    \[
        \mathrm{Per}_{\mathbb{R}^{n+1}}(\Omega_C', B_1) \leq \mathrm{Per}_{\mathbb{R}^{n+1}}(\Omega_C, B_1) - \delta
    \]
    for some $\delta = \delta(\theta) > 0$. As in~\cite[Lemma~3.4]{CLS22}, we choose $C^{1,\alpha}$-coordinates on $M$ around $Z$ with $g_{ij}(Z) = \delta_{ij}$ and rescale this competitor into $B_r(Z)$ to obtain a local area saving of $\delta(\theta) r^n (1 + o(1))$ as $r \to 0$. The same chord-beats-arc construction applies in the mass-cancellation case: the wedge geometry is identical regardless of whether $C$ arises as the boundary of a Caccioppoli set.

    \begin{figure}[ht]
        \centering
        \begin{tikzpicture}[scale=1.0]
          % Edge L: shared edge of the two parallelograms, with perspective
          \coordinate (A) at (-1.8,-0.6);   % near end of edge
          \coordinate (B) at (1.8,0.6);     % far end of edge

          % H_1 (upper half-plane): parallelogram above L
          \coordinate (A1) at ($(A)+(1.2,2.4)$);
          \coordinate (B1) at ($(B)+(1.2,2.4)$);

          % H_2 (lower half-plane): parallelogram below L
          \coordinate (A2) at ($(A)+(2.4,-1.2)$);
          \coordinate (B2) at ($(B)+(2.4,-1.2)$);

          % Fill both half-planes
          \fill[fillA, opacity=0.35] (A) -- (B) -- (B1) -- (A1) -- cycle;
          \fill[fillA, opacity=0.35] (A) -- (B) -- (B2) -- (A2) -- cycle;

          % Draw H_1 edges
          \draw[thick, bindA] (A) -- (A1);
          \draw[thick, bindA] (B) -- (B1);
          \draw[bindA] (A1) -- (B1);

          % Draw H_2 edges
          \draw[thick, bindA] (A) -- (A2);
          \draw[thick, bindA] (B) -- (B2);
          \draw[bindA] (A2) -- (B2);

          % Edge L (drawn on top)
          \draw[very thick] (A) -- (B);

          % Midpoint of L and ball B_r
          \coordinate (M) at ($(A)!0.5!(B)$);
          % Circle split at tangent points: short arc (right) solid, long arc (left) dashed
          \draw[gray, thin] ($(M)+(-18:1.2)$) arc[start angle=-18, end angle=55, radius=1.2];
          \draw[gray, thin, dashed] ($(M)+(55:1.2)$) arc[start angle=55, end angle=342, radius=1.2];
          \fill[gray] (M) circle (1.5pt);
          \node[gray, above left] at ($(M)+(135:1.2)$) {$\partial B_r$};

          % Two points on L inside the circle, on either side of M
          \coordinate (P) at ($(A)!0.35!(B)$);
          \coordinate (Q) at ($(A)!0.65!(B)$);
          \fill (P) circle (1.5pt);
          \fill (Q) circle (1.5pt);

          % Blue gradient shading for chord competitor (between semi-ellipses and chord)
          \shade[left color=blue!30, right color=blue!5, shading angle=18, opacity=0.5]
            (P) .. controls (-0.11, 0.69) and (0.50, 1.14) .. (0.69, 0.99)
            -- (1.14, -0.37)
            .. controls (1.08, -0.61) and (0.33, -0.61) .. cycle;

          % Gray shading for the wedge region (between chord and arc)
          \fill[gray, opacity=0.25]
            (0.69, 0.99) -- (1.14, -0.37)
            arc[start angle=-18, end angle=55, radius=1.2]
            -- cycle;

          % Semi-ellipse on H_1 side: solid left part, dashed right part
          \draw[thick] (P) .. controls (-0.11, 0.69) and (0.50, 1.14) .. (0.69, 0.99);
          \draw[thin, dashed] (0.69, 0.99) .. controls (0.79, 0.91) and (0.77, 0.65) .. (Q);

          % Semi-ellipse on H_2 side: solid left part, dashed right part
          \draw[thick] (P) .. controls (0.33, -0.61) and (1.08, -0.61) .. (1.14, -0.37);
          \draw[thin, dashed] (1.14, -0.37) .. controls (1.18, -0.25) and (1.01, -0.05) .. (Q);

          % Chord competitor: sparse parallel blue lines clipped to competitor region
          \begin{scope}
            \clip (P) .. controls (-0.11, 0.69) and (0.50, 1.14) .. (0.69, 0.99)
              -- (1.14, -0.37)
              .. controls (1.08, -0.61) and (0.33, -0.61) .. cycle;
            \foreach \s in {-0.15, 0.1, 0.35, 0.6} {
              \draw[blue, thin] ({\s-0.9}, {\s/3+2.72}) -- ({\s+0.9}, {\s/3-2.72});
            }
          \end{scope}

          % Rightmost solid blue line: shifted slightly right of T1-T2 to cover dashed ellipses
          \draw[blue, thin] (0.72, 0.97) -- (1.16, -0.37);

          % Past tangent points: dense dashed blue lines in the wedge region
          \begin{scope}
            \clip ($(M)+(-18:1.2)$)
              arc[start angle=-18, end angle=55, radius=1.2]
              .. controls (0.79, 0.91) and (0.77, 0.65) .. (Q)
              .. controls (1.01, -0.05) and (1.18, -0.25) .. cycle;
            \foreach \s in {-0.45,-0.35,...,0.85} {
              \draw[blue, thin, dashed] ({\s-0.9}, {\s/3+2.72}) -- ({\s+0.9}, {\s/3-2.72});
            }
          \end{scope}

          % Angle arc at far end B (interior angle, the Omega_C side)
          \draw[thin] (B) ++(-27:0.7) arc[start angle=-27, end angle=63, radius=0.7];
          \node at ($(B)+(18:1.0)$) {$\theta$};

          % Labels
          \node[bindA, above left] at ($(A1)!0.5!(B1)$) {$H_1$};
          \node[bindA, below right] at ($(A2)!0.5!(B2)$) {$H_2$};
          \node[above left] at ($(A)!0.3!(B)$) {$L$};
        \end{tikzpicture}
        \caption{The wedge (Case~1a): two half-hyperplanes $H_1, H_2$ meeting along the $(n-1)$-dimensional edge $L$ at angle $\theta < \pi$. In each two-dimensional cross-section normal to $L$, the two half-lines of $H_1 \cap \partial B_r$ and $H_2 \cap \partial B_r$ (black solid arcs) have total length $2r$, while the chord connecting their endpoints (blue lines) has length $2r\sin(\theta/2) < 2r$. Integrating along $L$, the blue surface has strictly less area than $\partial\Omega_C \cap B_r$ (grey region), so the wedge is not perimeter-minimizing.}
        \label{fig:wedge}
    \end{figure}

    \item \textbf{Case 1b (fold):} The tangent cone is $C = 2|H|$, where $H$ is a single half-hyperplane with multiplicity~$2$.

    In the no-mass-cancellation case, $E_i$ is a thin fold region between two sheets of $\spt\|V\|$ that both converge to $H$ at $Z$. Define $\Omega_r = \Omega \setminus (E_i \cap B_r(Z))$. Then $\Omega_r \subset \Omega$ and $\Omega \triangle \Omega_r = E_i \cap B_r(Z) \subset B_r(Z)$. The perimeter change has two contributions: we lose the sheet of $\partial^*\Omega$ separating $E_i$ from $\Omega^c$ inside $B_r(Z)$, saving $\sim \tfrac{1}{2}\omega_n r^n$ of area; and we gain the cross-section $\partial B_r(Z) \cap E_i$ as new boundary. Since $E_i$ is a thin fold closing up at $Z$, $\mathcal{H}^n(\partial B_r(Z) \cap E_i) = o(r^n)$. Therefore
    \[
        \mathrm{Per}_g(\Omega_r, B_r(Z)) - \mathrm{Per}_g(\Omega, B_r(Z)) = -r^n \delta(1 + o(1)),
    \]
    where $\delta = \omega_n / 2$. In the mass-cancellation case, the same collapsing construction applies to the two sheets of $V$ converging to $H$ at $Z$, yielding an identical area saving of $(\omega_n/2) r^n$.

    \item \textbf{Case 1c (higher-order junction, mass cancellation only):} The tangent cone consists of $N \geq 3$ half-hyperplanes meeting along a common $(n-1)$-dimensional edge.

    This case cannot arise when $V = |\partial^*\Omega(x_0)|$ (no mass cancellation), since the boundary of a Caccioppoli set separates $\mathbb{R}^{n+1}$ into two regions. In the mass-cancellation case, however, codimension-$1$ junctions with $N \geq 3$ sheets are possible. Such junctions are never area-minimizing in $\mathbb{R}^{n+1}$: an area-minimizing replacement in $B_r$ (which exists by~\cite[Theorem~5.1.6]{Fed69}) has strictly less mass than $C \mres B_r$, yielding a deficit $\delta r^n > 0$.
\end{itemize}

In every case, there exists $\delta > 0$ such that the junction structure at $Z$ admits a local area improvement of at least $\delta r^n$ in $B_r(Z)$.

\medskip\noindent\textbf{Step~2: Deriving a contradiction.}
The argument diverges depending on whether mass cancellation occurs.

\medskip\noindent\emph{Case A: No mass cancellation.}
Here $V = |\partial^*\Omega(x_0)|$. The area improvement from Step~1 shows that $\partial\Omega(x_0)$ is not one-sided homotopic minimizing in $B_r(Z)$ for all small $r$, so $Z \in \mathfrak{h}_{\mathrm{nm}}(V)$ (Definition~\ref{def:hnm}). By~\cite[Proposition~3.1]{CLS22}, $\mathfrak{h}_{\mathrm{nm}}(V)$ is finite, so there exists $\rho' > 0$ such that $\spt\|V\|$ is one-sided homotopic minimizing in $B_r(X)$ for every $X \in \spt\|V\| \setminus B_{\rho'}(Z)$ and some $r > 0$.

However, since the $\alpha$-structural hypothesis fails at $Z$, the common boundary $\Gamma \subset \sing V$ is $(n-1)$-dimensional. In particular, there exist singular points $X \in \sing V \setminus B_{\rho'}(Z)$ at which $V$ is one-sided homotopic minimizing. But $(n-1)$-dimensional singularities cannot occur for one-sided minimizers~\cite{Liu19}, a contradiction.

\medskip\noindent\emph{Case B: Mass cancellation.}
Here $V \neq |\partial^*\Omega(x_0)|$, so the one-sided minimizing argument of Case~A does not apply directly to $V$. Instead, we transfer the area improvement to the approximating sequence and derive a contradiction with the min-max width.

By varifold convergence $|\partial^*\Omega(x_i)| \to V$, for all $i$ sufficiently large $\partial^*\Omega(x_i) \cap B_r(Z)$ inherits a near-junction structure. Since each $\partial^*\Omega(x_i)$ is the smooth boundary of a Caccioppoli set (Proposition~\ref{prop:p2-federer-regularity}), the local competitor from Step~1 transfers via Caccioppoli surgery on $\Omega(x_i)$:
\begin{itemize}
    \item \emph{Wedge:} Set $\Omega(x_i)_r = \Omega(x_i) \setminus (W_i(r) \cap B_r(Z))$, excising the wedge tip.
    \item \emph{Fold:} Set $\Omega(x_i)_r = \Omega(x_i) \setminus (F_i(r) \cap B_r(Z))$, removing the thin fold layer.
    \item \emph{Higher-order junction:} Replace $\partial^*\Omega(x_i) \cap B_r(Z)$ by the area-minimizing spanning surface.
\end{itemize}
In each case, the perimeter decreases uniformly:
\[
    \mathrm{Per}(\Omega(x_i)) - \mathrm{Per}(\Omega(x_i)_r) \geq \frac{\delta r^n}{2} > 0
\]
for all $i$ large, where $\delta$ is independent of $i$.

Since $x_i \to x_0 \in \mathfrak{m}(\Phi)$, the perimeters satisfy $\mathrm{Per}(\Omega(x_i)) \to W$. The surgery modifies only the slices with parameters near $x_0$ (and is the identity elsewhere), producing a sweepout $\tilde{\Phi}$ with
\[
    \sup_x \mathrm{Per}(\tilde{\Omega}(x)) < W.
\]
But $W = \inf_{\Psi} \sup_x \mathrm{Per}(\Psi(x))$ is the min-max width, where the infimum is over \emph{all} sweepouts in the same homotopy class---not only nested ones. This contradicts the definition of $W$.

\medskip
In both cases, we reach a contradiction. Therefore, the $\alpha$-structural hypothesis holds.
\end{proof}

%%%%%%%%%%%%%%%%%%%%%%%%%%%%%%%%%%%%%%%%%%%%%%%%%%%%%%%%%%%%%%%%%%%%%%%%%%%%%%%
\subsection{Conclusion}\label{subsec:regularity-conclusion}
%%%%%%%%%%%%%%%%%%%%%%%%%%%%%%%%%%%%%%%%%%%%%%%%%%%%%%%%%%%%%%%%%%%%%%%%%%%%%%%

\begin{proof}[Proof of Theorem~\ref{thm:non-excessive-main}]
Let $(M^{n+1}, g)$ be a closed Riemannian manifold with $2 \leq n \leq 6$. By Theorem~\ref{thm:non-excessive-existence}, there exists a non-excessive ONVP sweepout $\Phi \in \Snm$. By Proposition~\ref{thm:CLS-stationary}, there exists a stationary varifold $V$ with $\mathbf{M}(V) = W$.

\medskip\noindent\emph{Case 1: No mass cancellation.}
By Theorem~\ref{thm:integrality-mult-one}, $V = |\partial^*\Omega(x_0)|$ is an integral varifold. By Proposition~\ref{prop:stability}, $V$ is stable. By Theorem~\ref{thm:alpha-hypo}, $V$ satisfies the $\alpha$-structural hypothesis. Therefore, all hypotheses of Wickramasekera's theorem (Theorem~\ref{thm:wickramasekera}) are satisfied, and $\sing V = \emptyset$.

\medskip\noindent\emph{Case 2: Mass cancellation.}
By Proposition~\ref{prop:integrality-nested}, $V$ is an integral varifold. By Proposition~\ref{prop:stability}, $V$ is stable. By Theorem~\ref{thm:alpha-hypo}, $V$ satisfies the $\alpha$-structural hypothesis. Again, all hypotheses of Wickramasekera's theorem are satisfied, and $\sing V = \emptyset$.

\medskip
In both cases, $\spt\|V\|$ is a smooth, closed, embedded minimal hypersurface in $M$.
\end{proof}
